\section{Benchmark design and evaluation protocol}
\label{sec:design}

\subsection{Data generation and leakage control}

Reads were simulated from a unified catalogue of human-gut reference genomes
\citep{Almeida2019} spanning 120 genera and 1{,}535 species, following the
read-simulation protocol of MetaTransformer \citep{Wichmann2023} with the ART
Illumina simulator \citep{Huang2012} at 150\,bp. The natural abundance
distribution is heavy-tailed (the most abundant genus, \emph{Collinsella},
accounts for 22.8\% of reads), which we retain to reflect realistic conditions
rather than an artificially uniform benchmark.

A subtlety specific to short-read simulation is that reads are not independent:
the source pool of 258.67M reads contains only 42.64M \emph{unique} reads
(6.07$\times$ redundancy), because high-coverage genomes generate overlapping
fragments. This matters for any data-scaling claim: past a point, adding
``more reads'' adds mostly repeats, not new information. We therefore report
scaling against the nominal training-read count while noting the underlying
unique-read count (50M reads $\rightarrow$ 17.8M unique; 250M $\rightarrow$
41.9M unique, near the 42.6M source ceiling).

All accuracy is reported on a held-out test pool that is strictly disjoint
from training at the read level. We use two test pools: a
\emph{natural-distribution} pool (99{,}742 reads, our primary metric) and a
harder \emph{leftover} pool as a conservative lower bound. For the Kraken\,2
comparison we additionally restrict to a database-coverage-matched pool of
85{,}819 reads whose source species are present in the Kraken\,2 database, so
that neural models and Kraken\,2 are scored on the same reads
(Section~\ref{sec:sample}).

\subsection{Models and the one-factor-at-a-time protocol}

The benchmark is organized so that each comparison isolates a single variable
(Figure~\ref{fig:design}):
\begin{itemize}\setlength{\itemsep}{2pt}
  \item \textbf{Data scale} --- the pre-trained NT-v2 backbone (498M
        parameters) fine-tuned with LoRA \citep{Hu2022} and an
        attention-pooling classification head, trained on 500K, 5M, 50M, and
        250M reads with everything else fixed.
  \item \textbf{Pre-training} --- the 29-layer pre-trained NT-v2 backbone
        versus a single-layer randomly initialized transformer, both with
        identical non-overlapping 6-mer tokenization and identical 50M-read
        data.
  \item \textbf{Tokenization} --- MetaTransformer-style encoders trained from
        scratch with 6-, 12-, and 13-mer tokenizers, overlapping (stride 1)
        versus non-overlapping (stride $=k$), on identical data.
  \item \textbf{Evaluation level} --- every model scored at the read level and,
        after aggregation, at the sample level, alongside Kraken\,2
        \citep{Wood2019}.
\end{itemize}
Because the axes are varied independently, the resulting deltas are more directly
attributable to each design factor rather than confounded by simultaneous changes
in model size, data, and tokenizer (this is a controlled ablation, not randomized
causal inference). Throughout, cross-tokenizer comparisons use the same nominal
read pools and the same held-out test pools, with unique-read counts reported
separately so that data-scaling claims are not confounded by read duplication.
Full hyperparameters, LoRA configuration, and training commands are given in the
Supplementary Material.

\begin{figure*}[t]
\centering
\resizebox{\textwidth}{!}{%
\begin{tikzpicture}[
    font=\footnotesize,
    node distance=6mm,
    box/.style={draw=bibblue,rounded corners=2pt,align=center,
                inner sep=4pt,fill=bibblue!5},
    axis/.style={draw=bibblue!70,rounded corners=2pt,align=center,
                 inner sep=4pt,fill=white},
    ev/.style={draw=bibblue,rounded corners=2pt,align=center,
               inner sep=4pt,fill=bibblue!12},
    ar/.style={-{Latex[length=2mm]},bibblue!80,thick}
  ]
  \node[box] (genomes) {Gut genome catalogue\\120 genera / 1{,}535 species};
  \node[box,right=of genomes] (reads) {ART 150\,bp reads\\(leakage-controlled\\train/test split)};
  \draw[ar] (genomes) -- (reads);

  % ablation axes
  \node[axis,above right=2mm and 10mm of reads] (a1) {\textbf{Data scale}\\500K\,$\to$\,250M};
  \node[axis,below=1.5mm of a1] (a2) {\textbf{Pre-training}\\pretrained vs scratch};
  \node[axis,below=1.5mm of a2] (a3) {\textbf{Tokenization}\\6/12/13-mer, overlap};
  \draw[ar] (reads.east) -- (a1.west);
  \draw[ar] (reads.east) -- (a2.west);
  \draw[ar] (reads.east) -- (a3.west);

  % evaluation
  \node[ev,right=14mm of a2] (e1) {\textbf{Read-level}\\Top-1 accuracy};
  \node[ev,above=1.5mm of e1] (e0) {\textbf{Sample-level}\\abundance ($r$, Bray--Curtis)};
  \node[ev,below=1.5mm of e1] (e2) {\textbf{Detection}\\ROC, sens.\ @ 95\% spec.};
  \draw[ar] (a1.east) -- (e0.west);
  \draw[ar] (a2.east) -- (e1.west);
  \draw[ar] (a3.east) -- (e2.west);

  \node[box,below=6mm of e1] (kraken) {Kraken\,2 baseline\\(coverage-matched pool)};
  \draw[ar] (e2.south) -- (kraken.north);
\end{tikzpicture}%
}
\caption{\textbf{Benchmark design.} Reads are simulated from a leakage-controlled
gut-genome catalogue. Three design axes (data scale, pre-training, tokenization)
are each varied one at a time while the others are held fixed, and every model
is scored at three evaluation levels---read-level accuracy, sample-level
abundance, and binary detection---against a coverage-matched Kraken\,2 baseline.}
\label{fig:design}
\end{figure*}
